% !TEX root = ../main.tex
\section{레지스트리와 \cpm}
\label{sec:registry}

\subsection{구성}

그림~\ref{fig:architecture}에 구성 요소를 보였다. 인덱서는 컨트랙트, 트랜잭션, 포스트컨디션, 자산 이벤트를 PostgreSQL 데이터베이스에 채운다. 과거 기록은 호스팅 서비스 대신 직접 운영하는 Stacks 노드와 직접 운영하는 Stacks Blockchain API에서 읽는다. 호스팅 서비스로 전체 이력을 채우면 요청 한도 때문에 느리고, 고정된 보관본이 있어야 평가를 재현할 수 있기 때문이다. 새 블록은 Chainhooks나 노드의 이벤트 인터페이스로 받는다 \cite{hiroapi,chainhooks2026}. 분석기는 컨트랙트를 하나씩 파싱해 \ref{sec:method}절의 참조와 자산 한도를 뽑고, 정규화 해시를 계산하고, 정적 검사를 돌린다. 매일 도는 작업이 그래프를 만들고 \PCER와 클론 묶음을 계산하고 검색 색인을 새로 고친다. 웹 화면과 \cpm\ 클라이언트는 같은 검색 API를 읽으며, \cpm은 받은 내용을 체인과 대조한다.

\begin{figure}[htbp]
\centering
\begin{tikzpicture}[
  box/.style={draw,rounded corners=2pt,align=center,font=\footnotesize,text width=2.5cm,minimum height=2.6em,inner sep=3pt},
  arr/.style={-{Stealth[length=2mm]},thick}]
\node[box] (node) at (0,0.4) {Stacks 노드와\\ Blockchain API\\ (보관본)};
\node[box] (hook) at (0,-2.2) {Chainhooks 또는\\ 이벤트\\ 인터페이스};
\node[box] (idx) at (3.2,-0.9) {인덱서};
\node[box] (db) at (6.4,-0.9) {PostgreSQL:\\ 컨트랙트,\\ 트랜잭션,\\ 포스트컨디션,\\ 이벤트};
\node[box] (ana) at (9.6,0.4) {분석기:\\ 참조, 자산 한도,\\ 해시, 검사};
\node[box] (rank) at (9.6,-2.2) {일일 작업:\\ 순위, 클론 묶음,\\ 검색 색인};
\node[box] (api) at (12.8,-2.2) {검색 API};
\node[box] (web) at (12.8,0.4) {웹 레지스트리};
\node[box] (cli) at (12.8,-4.0) {\cpm\ 클라이언트};
\draw[arr] (node) -- (idx);
\draw[arr] (hook) -- (idx);
\draw[arr] (idx) -- (db);
\draw[arr] (db) -- (ana);
\draw[arr] (db) -- (rank);
\draw[arr] (ana) -- (rank);
\draw[arr] (rank) -- (api);
\draw[arr] (api) -- (web);
\draw[arr] (api) -- (cli);
\end{tikzpicture}
\caption{Clarity Registry의 구성 요소. \cpm\ 클라이언트는 API가 돌려준 내용을 확인하려고 Stacks 노드에서 코드 해시를 직접 읽기도 한다.}
\label{fig:architecture}
\end{figure}

\subsection{패키지 모형}

\textbf{스코프.} 모든 컨트랙트에는 이미 게시자가 있다. principal 안의 주소가 그것이다. 그래서 레지스트리는 모든 컨트랙트를 그 principal 아래에 색인하고, 게시자는 읽기 쉬운 스코프를 두 가지 방법으로 차지할 수 있다. 주소 자체는 언제나 유효한 스코프이다. \code{alice.btc} 같은 BNS 이름은 차지를 신청할 때 그 이름이 배포자 주소로 해석되면 스코프가 된다. BNS 이름은 SIP-009 토큰이어서 팔 수 있으므로 \cite{bnsv2}, 레지스트리는 이름, 그 소유자, 신청한 블록 높이를 기록한다. 나중에 이름의 주인이 바뀌면 스코프는 동결된다. 기존 버전은 그대로 유지되고, 새 버전을 올리려면 새로 신청해야 하며, 스코프의 소유자가 프로젝트 잠금 파일에 기록된 소유자와 다르면 \cpm이 경고한다. 이것은 다른 레지스트리에서 알려진 공격인, 신뢰받는 이름의 탈취를 막는다 \cite{ladisa2023sok}.

\textbf{버전.} 배포된 컨트랙트는 바뀌지 않으므로 패키지의 버전마다 별개의 컨트랙트가 되고, 체인에서 \code{token-v1}과 \code{token-v2}는 서로 관계가 없다. 게시자는 그 관계를 매니페스트에 적는다. 패키지 이름, 각 컨트랙트의 시맨틱 버전, 각 버전의 이전 버전, 대응하는 테스트넷 principal, 라이선스 식별자, 저장소 URL, README, 감사 보고서 링크가 들어간다. 레지스트리는 나열된 모든 컨트랙트가 스코프의 주소로 배포되었고, 매니페스트가 그 주소의 SIP-018 구조화 데이터 서명을 갖추었거나 \cite{sip018} 그 주소에서 보낸 트랜잭션으로 제출된 경우에만 매니페스트를 받아들인다. 서명하는 주소가 곧 배포한 주소이므로 다른 키 기반 구조는 필요 없다.

\textbf{차지되지 않은 컨트랙트.} 매니페스트가 없는 컨트랙트도 principal 아래에 trait 표시, 클론 묶음, 점수와 함께 색인되지만, 버전과 README는 없다.

\subsection{명령}

\cpm은 Clarinet을 감싸는 도구이며 Clarinet을 대신하지 않는다. 표~\ref{tab:cpm}에 명령을 나열했다. 핵심은 참조와 복제의 구분이다. npm에서는 패키지를 설치하면 코드가 프로젝트로 복사된다. 프로젝트가 \emph{사용하려는} Clarity 컨트랙트는 이미 배포되어 모두가 함께 쓰며, principal로 호출된다. 프로젝트가 더하는 것은 참조이다. \code{cpm add}는 그 참조를 Clarinet 요구 사항으로 기록해서, 로컬 테스트는 Clarinet이 로컬 네트워크에 배포한 사본을 호출하고 메인넷 배포 때는 그 컨트랙트에 대해 아무것도 배포하지 않게 한다 \cite{clarinet}. SIP-010 템플릿으로 만든 토큰처럼 자기 인스턴스가 필요한 프로젝트는 \code{cpm create}나 \code{cpm fork}로 소스를 \emph{복제}해 자기 주소로 배포한다. 복제는 패키지의 라이선스가 허락할 때만 할 수 있다. 체인에 공개되었다고 오픈소스가 되는 것은 아니다.

\begin{table}[htbp]
\centering
\caption{\cpm\ 클라이언트의 명령}
\label{tab:cpm}
\small
\begin{tabularx}{\linewidth}{@{}>{\raggedright\arraybackslash}p{0.27\linewidth}>{\raggedright\arraybackslash}X@{}}
\toprule
명령 & 하는 일 \\
\midrule
\code{cpm search <query>} & 식~\eqref{eq:search}로 검색한다. trait, Clarity 버전, 네트워크로 거를 수 있다. \\
\code{cpm info <pkg>} & 버전, principal, 클론 묶음, 점수, 표~\ref{tab:signals}의 신호를 보여 준다. \\
\code{cpm add <pkg>[@ver]} & 배포된 컨트랙트를 참조한다. 그 컨트랙트와 정적 의존성을 찾고, 프로젝트와 Clarity 버전·에포크를 맞춰 보고, Clarinet 요구 사항을 더하고, 잠금 파일을 쓴다. \code{cpm install}은 별칭이다. \\
\code{cpm create <template> <name>} & 템플릿의 소스를 \code{contracts/}로 복사해 이름을 바꾸고, 선언된 매개변수를 채우고, 템플릿의 묶음 해시를 기록한다. 라이선스가 복제를 허락하지 않으면 거부한다. \\
\code{cpm fork <pkg>} & 라이선스가 허락하는 패키지에 대해 \code{create}와 같은 일을 한다. \\
\code{cpm why <contract>} & 프로젝트의 컨트랙트에서 주어진 컨트랙트까지의 의존 경로를 출력한다. \\
\code{cpm audit} & 모든 의존성에 대해 표~\ref{tab:signals}의 신호를 보고한다. \\
\code{cpm verify} & 잠긴 컨트랙트마다 노드에서 코드 해시를 읽어 잠금 파일과 비교한다. \\
\bottomrule
\end{tabularx}
\end{table}

\subsection{잠금 파일}

npm 잠금 파일은 \code{package.json}의 버전 범위가 실제로 풀린 정확한 버전을 고정한다. 배포된 컨트랙트에는 풀어야 할 버전이 없으므로, \cpm\ 잠금 파일은 신원을 기록한다. 의존성마다 메인넷 principal과 대응하는 테스트넷 principal, \code{contract-hash?}가 돌려주는 코드 해시, 배포 높이와 Clarity 버전, 묶음 해시, 라이선스, 추가할 때 본 스코프 소유자, 참조했는지 복제했는지를 적는다. \code{cpm verify}는 아래의 레지스트리 컨트랙트처럼 \code{contract-hash?}를 돌려주는 컨트랙트에 읽기 전용 호출을 보내 노드에서 코드 해시를 읽고, 잠금 파일과 비교한다. 그래서 API나 미러가 탈취되더라도 프로젝트가 의존하는 컨트랙트의 소스를 다른 것으로 바꿔치기할 수 없다.

\subsection{선택적인 레지스트리 컨트랙트}

이름을 체인에 둘 수도 있다. 작은 레지스트리 컨트랙트가 \code{publish(name, version, contract)}를 제공한다. 이 함수는 \code{principal-destruct?}로 \code{contract}의 주소 부분이 \code{tx-sender}와 같은지 확인하고, 그 컨트랙트의 \code{contract-hash?} 결과를 이름·버전과 함께 저장하고, 인덱서를 위해 이벤트를 출력한다. 아무 principal에 대해서나 \code{contract-hash?}를 돌려주는 읽기 전용 함수도 제공하며, \cpm은 이를 검증에 쓴다. 레지스트리는 이 컨트랙트 없이도 작동한다. 이 컨트랙트가 있으면 이름과 해시가 체인에 기록되어 웹 서비스 운영자에게 의존하지 않는다.

\subsection{사용자에게 보여 주는 신호}

순위는 결과의 순서를 정할 뿐, 컨트랙트가 안전하다고 말하지 않는다. 컨트랙트 페이지와 \code{cpm audit}은 점수 옆에 표~\ref{tab:signals}의 신호를 보여 주며, 레지스트리는 색인과 추천을 분리한다. 모든 컨트랙트를 색인한다. ``추천'' 표시를 둔다면 감사 링크, 최소 경과 기간, 심각한 지적 사항 없음처럼 점수 밖의 기준을 요구하며, 레지스트리는 점수만으로 컨트랙트를 악성이라 부르지 않는다.

\begin{table}[htbp]
\centering
\caption{컨트랙트마다 함께 보여 주는 신호}
\label{tab:signals}
\small
\begin{tabularx}{\linewidth}{@{}>{\raggedright\arraybackslash}p{0.34\linewidth}>{\raggedright\arraybackslash}X@{}}
\toprule
신호 & 중요한 이유 \\
\midrule
\code{with-all-assets-unsafe}, 또는 다른 컨트랙트 호출을 감싼 옛 \code{as-contract} & 컨트랙트가 다른 코드에 한도 없이 자기 자산을 옮기게 한다. STACY는 \code{as-contract} 안의 호출을 찾아낸다 \cite{stacy}. \\
\code{as-contract?} 안에서 호출하는 trait 타입 인자 & 컨트랙트의 자산이 호출자가 고른 코드에 노출된다. \\
Allow 모드 호출의 비율 & 이 컨트랙트의 사용자들이 포스트컨디션 보호 없이 트랜잭션에 서명하고 있다. \\
클론 묶음과 대표 & 이 컨트랙트가 묶음에서 가장 많이 쓰이는 구성원인지, 많은 복제본 가운데 하나인지 보여 준다. \\
경과 기간, Clarity 버전, 에포크 & 오래된 Clarity 버전에는 자산 한도가 없다. 아주 새로운 컨트랙트는 이력이 짧다. \\
정적 분석 지적 사항 & STACY와 레지스트리 자체 검사의 결과. \\
감사 링크 & 매니페스트에서 가져온다. 레지스트리는 감사 자체를 검증하지 않는다. \\
\bottomrule
\end{tabularx}
\end{table}
